\documentclass[reqno]{amsart}
\usepackage[T1]{fontenc}
\usepackage{amsmath,amssymb,amsthm}
\usepackage{xr-hyper}
\usepackage[hidelinks]{hyperref}
\externaldocument[sm-]{20261125T000000--slope-method__workbook}
\externaldocument[cv-]{20261125T006000--covolumes__workbook}
\externaldocument[cf-]{20261125T017500--chow-forms-and-effective-hilbert-samuel__workbook}
\externaldocument[hp-]{20261125T020300--heights-periods-and-the-auxiliary-prime__workbook}
\externaldocument[sg-]{20261126T020320--slopes-of-lattices-of-forms__workbook}
\externaldocument[fl-]{20261126T020330--the-cm-period-and-the-flat-line__workbook}

\newcounter{chap}
% Chapter numbers, generated from 20261125T012000--grh-workbook__guide.org
% by scripts/make-chapters.py.  Do not edit: change the order of the
% index and rerun the script (./preview-site.sh and CI do this).
%
% Usage in a sheet, after \newcounter{chap}:
%   % Chapter numbers, generated from 20261125T012000--grh-workbook__guide.org
% by scripts/make-chapters.py.  Do not edit: change the order of the
% index and rerun the script (./preview-site.sh and CI do this).
%
% Usage in a sheet, after \newcounter{chap}:
%   % Chapter numbers, generated from 20261125T012000--grh-workbook__guide.org
% by scripts/make-chapters.py.  Do not edit: change the order of the
% index and rerun the script (./preview-site.sh and CI do this).
%
% Usage in a sheet, after \newcounter{chap}:
%   \input{chapters}\chapnumber{<sheet ID>}
\makeatletter
\providecommand\chapters@set[2]{\@namedef{chapters@#1}{#2}}
\chapters@set{20261125T001000}{1}% zeros-to-ramp-sums
\chapters@set{20261125T001100}{2}% inert-primes-and-estimate-b
\chapters@set{20261125T005000}{3}% galois-fourier-and-averaging
\chapters@set{20261125T005500}{4}% columns-from-a-sphere
\chapters@set{20261125T000000}{5}% slope-method
\chapters@set{20261126T030100}{6}% a-twist-for-each-prime
\chapters@set{20261125T003000}{7}% the-chain-and-the-square-root
\chapters@set{20261125T010000}{8}% rank-lemma-and-zero-estimate
\chapters@set{20261125T010500}{9}% poincare-series-and-near-orthogonality
\chapters@set{20261124T223000}{10}% elliptic-curves-for-the-proof
\chapters@set{20261125T009500}{11}% fock-model-and-theta-intertwiners
\chapters@set{20261124T223100}{12}% complex-multiplication-and-inert-reduction
\chapters@set{20261125T008000}{13}% one-cm-determinant
\chapters@set{20261124T214500}{14}% supersingular-theta-divisibility
\chapters@set{20261124T214600}{15}% the-isotropic-plane
\chapters@set{20261125T013700}{16}% pairings-on-a-formal-group-of-height-two
\chapters@set{20261125T013500}{17}% hecke-nondegeneracy-on-the-tree
\chapters@set{20261126T030200}{18}% the-bound-is-attained
\chapters@set{20261125T009000}{19}% kronecker-limit-and-chowla-selberg
\chapters@set{20261125T002000}{20}% definite-theta-lifts
\chapters@set{20261125T015500}{21}% siegel-weil-in-the-smallest-case
\chapters@set{20261125T002100}{22}% definite-theta-lifts-and-norms
\chapters@set{20261125T016500}{23}% hecke-operators-and-theta-lifts
\chapters@set{20261125T016700}{24}% divisibility-along-the-chain
\chapters@set{20261125T004000}{25}% denominators-and-sha
\chapters@set{20261125T004100}{26}% tate-twists-and-sha
\chapters@set{20261126T020430}{27}% the-minus-part-by-gauss-sums
\chapters@set{20261125T014100}{28}% congruence-modules
\chapters@set{20261125T014500}{29}% eisenstein-congruences-and-selmer-classes
\chapters@set{20261125T014700}{30}% pseudocharacters-and-the-length-bound
\chapters@set{20261125T011000}{31}% galois-representations-black-box
\chapters@set{20261126T030300}{32}% the-data
\chapters@set{20261126T020210}{33}% bessel-poincare-series-and-their-gram-operators
\chapters@set{20261126T030400}{34}% the-seesaw-identity
\chapters@set{20261126T030500}{35}% divisibility-of-the-composite-lift
\chapters@set{20261126T030600}{36}% theorem-a-short-forms
\chapters@set{20261126T030700}{37}% the-determinant-of-a-chain
\chapters@set{20261126T030800}{38}% the-cost-of-one-lift
\chapters@set{20261126T030900}{39}% denominators-of-projections-to-lifts
\chapters@set{20261126T031000}{40}% twists-with-few-congruences
\chapters@set{20261126T031100}{41}% theorem-c-and-the-main-inequality
\providecommand\chapnumber[1]{%
  \@ifundefined{chapters@#1}%
    {\PackageError{chapters}{Sheet #1 is not listed in the index}%
      {List it in 20261125T012000--grh-workbook__guide.org\MessageBreak
       and run scripts/make-chapters.py.}}%
    {\edef\thechap{\@nameuse{chapters@#1}}}}
\makeatother
\chapnumber{<sheet ID>}
\makeatletter
\providecommand\chapters@set[2]{\@namedef{chapters@#1}{#2}}
\chapters@set{20261125T001000}{1}% zeros-to-ramp-sums
\chapters@set{20261125T001100}{2}% inert-primes-and-estimate-b
\chapters@set{20261125T005000}{3}% galois-fourier-and-averaging
\chapters@set{20261125T005500}{4}% columns-from-a-sphere
\chapters@set{20261125T000000}{5}% slope-method
\chapters@set{20261126T030100}{6}% a-twist-for-each-prime
\chapters@set{20261125T003000}{7}% the-chain-and-the-square-root
\chapters@set{20261125T010000}{8}% rank-lemma-and-zero-estimate
\chapters@set{20261125T010500}{9}% poincare-series-and-near-orthogonality
\chapters@set{20261124T223000}{10}% elliptic-curves-for-the-proof
\chapters@set{20261125T009500}{11}% fock-model-and-theta-intertwiners
\chapters@set{20261124T223100}{12}% complex-multiplication-and-inert-reduction
\chapters@set{20261125T008000}{13}% one-cm-determinant
\chapters@set{20261124T214500}{14}% supersingular-theta-divisibility
\chapters@set{20261124T214600}{15}% the-isotropic-plane
\chapters@set{20261125T013700}{16}% pairings-on-a-formal-group-of-height-two
\chapters@set{20261125T013500}{17}% hecke-nondegeneracy-on-the-tree
\chapters@set{20261126T030200}{18}% the-bound-is-attained
\chapters@set{20261125T009000}{19}% kronecker-limit-and-chowla-selberg
\chapters@set{20261125T002000}{20}% definite-theta-lifts
\chapters@set{20261125T015500}{21}% siegel-weil-in-the-smallest-case
\chapters@set{20261125T002100}{22}% definite-theta-lifts-and-norms
\chapters@set{20261125T016500}{23}% hecke-operators-and-theta-lifts
\chapters@set{20261125T016700}{24}% divisibility-along-the-chain
\chapters@set{20261125T004000}{25}% denominators-and-sha
\chapters@set{20261125T004100}{26}% tate-twists-and-sha
\chapters@set{20261126T020430}{27}% the-minus-part-by-gauss-sums
\chapters@set{20261125T014100}{28}% congruence-modules
\chapters@set{20261125T014500}{29}% eisenstein-congruences-and-selmer-classes
\chapters@set{20261125T014700}{30}% pseudocharacters-and-the-length-bound
\chapters@set{20261125T011000}{31}% galois-representations-black-box
\chapters@set{20261126T030300}{32}% the-data
\chapters@set{20261126T020210}{33}% bessel-poincare-series-and-their-gram-operators
\chapters@set{20261126T030400}{34}% the-seesaw-identity
\chapters@set{20261126T030500}{35}% divisibility-of-the-composite-lift
\chapters@set{20261126T030600}{36}% theorem-a-short-forms
\chapters@set{20261126T030700}{37}% the-determinant-of-a-chain
\chapters@set{20261126T030800}{38}% the-cost-of-one-lift
\chapters@set{20261126T030900}{39}% denominators-of-projections-to-lifts
\chapters@set{20261126T031000}{40}% twists-with-few-congruences
\chapters@set{20261126T031100}{41}% theorem-c-and-the-main-inequality
\providecommand\chapnumber[1]{%
  \@ifundefined{chapters@#1}%
    {\PackageError{chapters}{Sheet #1 is not listed in the index}%
      {List it in 20261125T012000--grh-workbook__guide.org\MessageBreak
       and run scripts/make-chapters.py.}}%
    {\edef\thechap{\@nameuse{chapters@#1}}}}
\makeatother
\chapnumber{<sheet ID>}
\makeatletter
\providecommand\chapters@set[2]{\@namedef{chapters@#1}{#2}}
\chapters@set{20261125T001000}{1}% zeros-to-ramp-sums
\chapters@set{20261125T001100}{2}% inert-primes-and-estimate-b
\chapters@set{20261125T005000}{3}% galois-fourier-and-averaging
\chapters@set{20261125T005500}{4}% columns-from-a-sphere
\chapters@set{20261125T000000}{5}% slope-method
\chapters@set{20261126T030100}{6}% a-twist-for-each-prime
\chapters@set{20261125T003000}{7}% the-chain-and-the-square-root
\chapters@set{20261125T010000}{8}% rank-lemma-and-zero-estimate
\chapters@set{20261125T010500}{9}% poincare-series-and-near-orthogonality
\chapters@set{20261124T223000}{10}% elliptic-curves-for-the-proof
\chapters@set{20261125T009500}{11}% fock-model-and-theta-intertwiners
\chapters@set{20261124T223100}{12}% complex-multiplication-and-inert-reduction
\chapters@set{20261125T008000}{13}% one-cm-determinant
\chapters@set{20261124T214500}{14}% supersingular-theta-divisibility
\chapters@set{20261124T214600}{15}% the-isotropic-plane
\chapters@set{20261125T013700}{16}% pairings-on-a-formal-group-of-height-two
\chapters@set{20261125T013500}{17}% hecke-nondegeneracy-on-the-tree
\chapters@set{20261126T030200}{18}% the-bound-is-attained
\chapters@set{20261125T009000}{19}% kronecker-limit-and-chowla-selberg
\chapters@set{20261125T002000}{20}% definite-theta-lifts
\chapters@set{20261125T015500}{21}% siegel-weil-in-the-smallest-case
\chapters@set{20261125T002100}{22}% definite-theta-lifts-and-norms
\chapters@set{20261125T016500}{23}% hecke-operators-and-theta-lifts
\chapters@set{20261125T016700}{24}% divisibility-along-the-chain
\chapters@set{20261125T004000}{25}% denominators-and-sha
\chapters@set{20261125T004100}{26}% tate-twists-and-sha
\chapters@set{20261126T020430}{27}% the-minus-part-by-gauss-sums
\chapters@set{20261125T014100}{28}% congruence-modules
\chapters@set{20261125T014500}{29}% eisenstein-congruences-and-selmer-classes
\chapters@set{20261125T014700}{30}% pseudocharacters-and-the-length-bound
\chapters@set{20261125T011000}{31}% galois-representations-black-box
\chapters@set{20261126T030300}{32}% the-data
\chapters@set{20261126T020210}{33}% bessel-poincare-series-and-their-gram-operators
\chapters@set{20261126T030400}{34}% the-seesaw-identity
\chapters@set{20261126T030500}{35}% divisibility-of-the-composite-lift
\chapters@set{20261126T030600}{36}% theorem-a-short-forms
\chapters@set{20261126T030700}{37}% the-determinant-of-a-chain
\chapters@set{20261126T030800}{38}% the-cost-of-one-lift
\chapters@set{20261126T030900}{39}% denominators-of-projections-to-lifts
\chapters@set{20261126T031000}{40}% twists-with-few-congruences
\chapters@set{20261126T031100}{41}% theorem-c-and-the-main-inequality
\providecommand\chapnumber[1]{%
  \@ifundefined{chapters@#1}%
    {\PackageError{chapters}{Sheet #1 is not listed in the index}%
      {List it in 20261125T012000--grh-workbook__guide.org\MessageBreak
       and run scripts/make-chapters.py.}}%
    {\edef\thechap{\@nameuse{chapters@#1}}}}
\makeatother

\ifcsname chapters@20261126T020340\endcsname
  \chapnumber{20261126T020340}
\else
  \renewcommand{\thechap}{?}
\fi
\renewcommand{\thesection}{\thechap.\arabic{section}}
\theoremstyle{definition}
\newtheorem{problem}{Problem}[section]
\theoremstyle{plain}
\newtheorem{theorem}{Theorem}[section]
\theoremstyle{remark}
\newtheorem*{leanremark}{In the Lean formalization}
\newcommand{\Z}{\mathbb Z}
\newcommand{\Q}{\mathbb Q}
\newcommand{\C}{\mathbb C}
\newcommand{\Pb}{\mathbb P}
\newcommand{\Oc}{\mathcal O}
\newcommand{\rk}{\operatorname{rk}}
\newcommand{\mus}{\mu^\sharp}
\newcommand{\dg}{\widehat{\deg}}
\newcommand{\cc}{\hat{c}_1}
\newcommand{\vol}{\widehat{\mathrm{vol}}}
\newcommand{\rd}{\operatorname{rd}}
\newcommand{\tl}[1]{\texttt{#1}}
\setlength{\emergencystretch}{3em}

\title{Lower slope bounds on covers}
\author{}
\date{November 26, 2026}

\begin{document}
\maketitle
\refstepcounter{chap}\label{chap}

\paragraph{Prerequisites.}
From Chapter~\ref{sm-chap}: \(\mus\) and Problem~\ref{sm-pr:additivity}.  From Chapter~\ref{cv-chap}, \emph{Covolumes of lattices of sections}: the lattices \(E_D\), Problems~\ref{cv-pr:degED} and~\ref{cv-pr:degED-asym}, Theorem~\ref{cv-thm:HS} and formula~\eqref{cv-eq:hPN}.  From Chapter~\ref{sg-chap}: Lemma~\ref{sg-lem:descent} and the notation of Section~\ref{sg-sec:reduction} (the lift \(j\), \(y = x^\varepsilon\), \(b = e^y\), the error unit \(u = B_0\log\log x \ge 1\), the main term \(T_j\), and \(\mu_{\mathrm{core}}\)).  From Chapter~\ref{fl-chap}: Theorem~\ref{fl-thm:flatat}, and the notation \(d\), \(c_\chi\), \(E\), \(\rd_E\) of Section~\ref{fl-sec:constant}.  Heights of points and arithmetic intersection numbers are recalled in Section~\ref{sec:heights} with references.

Chapter~\ref{sg-chap} reduced the covolume estimate to a lower bound \(\mu_{\mathrm{core}} \ge -T_j - C_Lu\) for the slope of a lattice of holomorphic cusp forms on the Shimura variety \(S_j\) of the \(j\)th lift.  How can one bound a slope from below with an error that is a fixed multiple of \(u\), uniformly in the point and the lift?  In the Lean formalization \cite{Lean}, for \(j < m\), one passes to a cover of \(S_j\), changes the weight, applies an effective arithmetic Hilbert--Samuel inequality on each component, and descends.  The last lift, \(j = m\), has \(q_m = 0\) and is treated separately; we have not read that proof.

Throughout, \(F\) is a number field, the field of definition of the cover; for a lattice \(\Lambda\) of sections we write \(\dg\Lambda := \deg\Lambda/[F:\Q]\), so that \(\mus(\Lambda) = 2\dg\Lambda/\rk\Lambda\).

\section{Hilbert--Samuel on one component}\label{sec:hs}

Let \(Z\) be an integral projective arithmetic variety over \(\Oc_F\) whose generic fibre has dimension \(\nu\), and \(L\) a line bundle on \(Z\) with a smooth Hermitian metric invariant under complex conjugation, ample on the generic fibre.  Write \(\delta_Z := \deg L_F\) and, following Zhang \cite{Zha},
\[
  h_Z := \frac{\cc(L)^{\nu+1}}{(\nu+1)\,[F:\Q]\,\delta_Z} .
\]
For an integer \(T \ge 1\) give \(H^0(Z, L^{\otimes T})\) its \(L^2\) norms; we use the inequality
\begin{equation}\label{eq:hs}
  \dg H^0(Z, L^{\otimes T}) \ge \binom{T-1}{\nu}\,\delta_Z\,\bigl(Th_Z - 3(\nu+1)\lvert h_Z\rvert - A\log(2 + \lvert h_Z\rvert) - R\bigr) ,
\end{equation}
where \(A > 0\) and \(R \ge 0\) are numbers with \(A\log\max(A,1) + R \le C_2u\).

\emph{Status}: in the formalization~\eqref{eq:hs} is proved, for the line bundles and components that occur and for the tensor power \(T\) chosen there, from published theorems: Kawamata--Viehweg and Nadel vanishing, principalisation, toroidal compactification, and the arithmetic Hilbert--Samuel theorem; the proof also uses Chow forms (Chapter~\ref{cf-chap}).  We have read its statement, not its proof, and we have not checked that the formalization's height is normalized as \(h_Z\) above.  Problem~\ref{pr:hs-leading} shows that the leading terms agree with this normalization.  Unlike Theorem~\ref{cv-thm:HS}, \eqref{eq:hs} is effective: it holds at one \(T\), with all errors explicit.

\begin{problem}\label{pr:binom}
For \(T > \nu \ge 1\), \((T-\nu)^\nu/\nu! \le \binom{T-1}{\nu} \le T^\nu/\nu!\).  For \(T = 100\) and \(\nu = 3\) the three numbers are \(152112.1\ldots\), \(156849\) and \(166666.6\ldots\)
\end{problem}

\begin{problem}\label{pr:hs-leading}
For fixed \(Z\), \(L\), \(A\) and \(R\), the right side of~\eqref{eq:hs} is \(\cc(L)^{\nu+1}T^{\nu+1}/((\nu+1)!\,[F:\Q]) + O(T^\nu)\) as \(T \to \infty\), the leading term of Theorem~\ref{cv-thm:HS}.
\end{problem}

\begin{problem}\label{pr:hs-P1}
For \(Z = \Pb^1_\Z\) and \(\Oc(1)\) with the Fubini--Study metric, \(\nu = 1\), \(\delta_Z = 1\) and \(h_Z = \frac14\).  With \(A\) and \(R\) replaced by \(0\), \eqref{eq:hs} reads \(\deg E_T \ge (T-1)(\frac T4 - \frac32)\).  This holds for \(1 \le T \le 400\), as one checks by machine, and for all large \(T\) by Problem~\ref{cv-pr:degED-asym}; for instance \(\deg E_{10} = 32.1969\ldots \ge 9\) and \(\deg E_{50} = 677.4825\ldots \ge 539\).
\end{problem}

\begin{problem}\label{pr:log-term}
For \(y \ge 0\) and \(A > 0\), \(2y - A\log(2 + y) \ge -A\log\max(A,1) - 1\).
\end{problem}

\begin{problem}\label{pr:abs-trick}
Let \(c_s \ge 0\), \(h \ge -c_s\) and \(T \ge 3\nu + 5\).  Then
\[
  Th - 3(\nu+1)\lvert h\rvert - A\log(2 + \lvert h\rvert) \ge (T - 3\nu - 5)\,h - (6\nu + 10)\,c_s - A\log\max(A,1) - 1 .
\]
\end{problem}

\section{From the components to the cover}\label{sec:cover}

Let \(Y\) be a cover of \(S_j\), \(j < m\), of the kind described in Section~\ref{sec:descend}.  In the formalization a weight \(\lambda_M\) is fixed which is \(T\) times a weight \(\lambda_1\), and the integral cusp forms on \(Y\) of weight \(\lambda_M\) are compared with the sections of \(L^{\otimes T}\) on the components \(Z\) of a model attached to \(Y\) and \(\lambda_1\).  The number \(\nu\) is attached to the signature; \(T\) is chosen of size about \(\nu^2\) times a large parameter, so that \(T \ge 3\nu + 5\).  Write \(\mu_Y(\lambda) := \mus\) of the lattice of integral cusp forms of weight \(\lambda\) on \(Y\), and, for each component \(Z\), \(r_Z\) and \(\dg_Z\) for the rank and the normalized degree of its lattice of sections.  Let \(B := \binom{T-1}{\nu}\).  The formalization proves:
\begin{itemize}
\item[(U1)] \(\mu_Y(\lambda_M) \ge 2\sum_Z\dg_Z\big/\sum_Zr_Z\);
\item[(U2)] \(r_Z \ge B\delta_Z\) for every \(Z\);
\item[(U3)] there are \(\delta' > 0\) and \(0 \le \varrho < 1\) with \(\delta' \le \delta_Z \le \delta'/(1-\varrho)\) for every \(Z\);
\item[(U4)] there are \(c_s, c_c, X \ge 0\) and real numbers \(s_Z\) with \(h_Z + c_s \ge 0\), \((1-\varrho)\max(0, s_Z + c_s) \le h_Z + c_s + c_c\), and \(\frac1N\sum_Zs_Z \ge -X\), where \(N\) is the number of components;
\item[(U5)] the inequality~\eqref{eq:hs} on every \(Z\);
\item[(U6)] \(TX \le C_1u\), \(Tc_c \le C_1u\) and \((6\nu + 10 + 2\varrho T)\,c_s \le C_1u\).
\end{itemize}

\begin{problem}\label{pr:sum-orth}
If the lattice of forms on \(Y\) is the orthogonal sum of the lattices of sections on the components, then (U1) is an equality.
\end{problem}

\begin{problem}\label{pr:weighted-mean}
Under (U3) and (U4),
\[
  \frac{\sum_Z\delta_Zh_Z}{\sum_Z\delta_Z} \ge -X - 2\varrho c_s - c_c .
\]
\end{problem}

\begin{problem}\label{pr:weighted-example}
Let \(\varrho = \frac12\), two components with \(\delta = (1, 2)\), \(s = (2, -2)\), \(c_s = 1\), \(c_c = 0\), \(X = 0\), and \(h = (\frac12, -1)\).  Then (U3) and (U4) hold, the weighted mean of \(h\) is \(-\frac12\), its plain mean is \(-\frac14\), and the bound of Problem~\ref{pr:weighted-mean} is \(-1\).
\end{problem}

\begin{problem}\label{pr:core}
Under (U1)--(U6),
\begin{align*}
  \mu_Y(\lambda_M) &\ge -2\bigl(TX + Tc_c + (6\nu + 10 + 2\varrho T)c_s + A\log\max(A,1) + R + 1\bigr)\\
  &\ge -2(3C_1 + C_2 + 1)\,u .
\end{align*}
\end{problem}

Problems~\ref{pr:abs-trick}, \ref{pr:weighted-mean} and~\ref{pr:core} prove the following.

\begin{theorem}[uniform lower bound]\label{thm:uniform}
Under (U1)--(U6), \(\mu_Y(\lambda_M) \ge -C_Uu\) with \(C_U = 2(3C_1 + C_2 + 1)\).
\end{theorem}

The formalization proves Theorem~\ref{thm:uniform} with \(4\) in place of \(1\) in \(C_U\), under a version of (U2) with an upper bound for \(r_Z\) as well, which the argument above does not need.  The content is in (U4) and (U6): the heights of the components must be bounded below on average by \(-X\) with \(TX \le C_1u\).

\section{Heights of components}\label{sec:heights}

Where does (U4) come from?  For an algebraic point \(P\) of \(Z_F\) over a finite extension \(F'\), extended to \(\Oc_{F'}\), let \(h_L(P) := \deg(P^*L)/[F':\Q]\).  For a real \(c\), let \(L(-c)\) be \(L\) with its norms multiplied by \(e^c\); then \(h_{L(-c)}(Z) = h_Z - c\) \cite{Zha}.

\begin{problem}\label{pr:rescale-height}
\(h_{L(-c)}(P) = h_L(P) - c\) for every algebraic point \(P\).
\end{problem}

\begin{problem}\label{pr:small-section}
If \(s \in H^0(Z, L^{\otimes k})\) is integral with \(\lVert s\rVert_{\sup} \le e^{-kc}\) at every archimedean place, and \(P\) is not in the zero locus of \(s\), then \(h_L(P) \ge c\).
\end{problem}

Zhang used Problem~\ref{pr:small-section} and Minkowski's theorem in the proof of his inequalities for successive minima \cite[proof of Thm.~5.2]{Zha}.  Together with Moriwaki's arithmetic Hilbert--Samuel formula for nef bundles \cite{Mo} the same argument gives two comparisons, as explained in \cite[\S4.3]{Zh}: if \(h_L(P) \le b\) on a Zariski dense set of points, then \(h_Z \le b\); and if \(h_L(P) \ge a\) for all points, then \(h_Z \ge a\).

\begin{problem}\label{pr:P1-points}
On \(\Pb^1_\Z\) with the Fubini--Study metric, \(h((p:q)) = \frac12\log(p^2 + q^2)\) for coprime integers \(p, q\), and \(h((1:\zeta)) = \frac12\log2 = 0.3465\ldots\) for every root of unity \(\zeta\).  The second family is Zariski dense, which is consistent with \(h_Z = \frac14\).
\end{problem}

Zariski dense points of small height also exist on Grassmannians: for an \(a \times a'\) matrix \(W\), the row space of \((I_a\ W)\) is a point of the Grassmannian of \(a\)-planes in \(F^{a+a'}\), and \(\frac12\log\det(I_a + WW^*)\) is the archimedean part of its Pl\"ucker height for the Fubini--Study metric \cite[\S4.3]{Zh}.

\begin{problem}\label{pr:grassmann}
If \(W\) is an \(a \times a'\) matrix whose entries are roots of unity, then \(\frac12\log\det(I_a + WW^*) \le \frac a2\log(a' + 1)\).  For \(a = a' = 2\) and \(W = \bigl(\begin{smallmatrix}1&1\\1&-1\end{smallmatrix}\bigr)\) both sides equal \(\log3\).
\end{problem}

The lower bound for heights needs small sections when \(L\) is only nef.  Moriwaki proved that if \(L\) is nef on every fibre and on the generic fibre, with a smooth semipositive metric invariant under conjugation, then \(\vol(L) \ge \cc(L)^{\nu+1}\), where \(\vol\) is the arithmetic volume of Chapter~\ref{cv-chap} \cite{Mo}, after Yuan \cite{Yu}.  \emph{Status}: this statement is proved in the formalization, from resolution of singularities, a Poincar\'e--Lelong formula and a density statement for curves; we have not read below these.

At the finite places, the heights of points on the Shimura variety involve the integral structure of the invariant differentials, which the formalization controls with the following theorem.  Let \(C\) be an algebraically closed complete extension of \(\Q_p\), \(G\) a \(p\)-divisible group over \(\Oc_C\), and \(\alpha_G \colon T_p(G) \otimes \Oc_C \to \omega_{G^\vee}\) its Hodge--Tate map: an element \(x\) of \(T_p(G) = \operatorname{Hom}(G^\vee, \mu_{p^\infty})\) is sent to the pullback \(x^*(dT/T)\) of the invariant differential of \(\mu_{p^\infty}\).  Then \(\kappa\cdot\omega_{G^\vee} \subseteq \operatorname{im}\alpha_G\) for every \(\kappa\) with \(v_p(\kappa) \ge 1/(p-1)\) (with \(v_p(\kappa) \ge 2\) if \(p = 2\)) \cite{Fa}, and the kernel of \(\alpha_G\) is the annihilator, under the Tate pairing, of the kernel of the logarithmic differential of \(G^\vee\) \cite{ALB}.  \emph{Status}: the second statement is assumed in the formalization as a published theorem, and the first is derived from it there.

\begin{problem}\label{pr:HT-mult}
For \(G = \Q_p/\Z_p\), \(\alpha_G\) maps \(T_p(G) \otimes \Oc_C\) onto \(\omega_{\mu_{p^\infty}} = \Oc_C\,dT/T\); for \(G = \mu_{p^\infty}\), \(\omega_{G^\vee} = 0\).
\end{problem}

In the formalization, (U4) and (U6) follow from the following statement on covers.  For every lift \(j < m\) with \(m^2 \le k\), and every good cover \(Y\) of \(S_j\) whose level has at most fifteen prime factors, the mean over the components of certain volume slopes \(s_Z\) is at least \(-X\), where \(X\) is a constant times an error depending on the weight and on \(\log\rd_E + \lvert c_\chi\rvert\).  \emph{Status}: proved in the formalization from Moriwaki's theorem, the arithmetic Hilbert--Samuel formula for nef bundles, an absolute Siegel lemma \cite[\S4.3]{Zh}, resolution, the extension of homomorphisms of semi-abelian schemes, the Hodge--Tate theorem above, Zarhin's trick, and a Riemann--Roch statement for nef bundles.  We have not followed this proof, which goes through towers of covers and their volumes, and we do not state the error term.

\section{Back to the core level}\label{sec:descend}

A \emph{good cover} of \(S_j\) in the formalization has full level \(15^eg_0g_1\), with an exponent \(e\) fixed by the size parameters of the family and \(g_0, g_1\) the level primes, and a field of definition of bounded size; such covers exist by class field theory.  Besides \(\lambda_M\) it carries weights \(\lambda_-\) and \(\lambda_+\), shifts of the weight of the \(j\)th lift, and a lattice \(L_{\mathrm{deep}}\) of forms on \(Y\) pulled back from the core level.  With a size parameter \(\Lambda_j > 0\), the formalization proves:
\begin{itemize}
\item[(D1)] \(1 \le b \le u\), \(\lvert T_j\rvert \le \Lambda_j\) and \(b \le \Lambda_j\);
\item[(D2)] \(\mus(L_{\mathrm{deep}}) - 2b \le \mu_{\mathrm{core}}\);
\item[(D3)] there are \(n\) maps from the forms of weight \(\lambda_-\) to the forms of the weight of the \(j\)th lift on \(Y\), jointly injective, mapping integral forms into \(L_{\mathrm{deep}}\), with \(\sum_h\lVert\varphi_hf\rVert^2 \le e^{2b}\lVert f\rVert^2\), and \(n\rk L_{\mathrm{deep}}\) is at most the rank for \(\lambda_+\);
\item[(D4)] the rank for \(\lambda_+\) is at most \((1 + b/\Lambda_j)\) times the rank for \(\lambda_-\);
\item[(D5)] \(L_{\mathrm{deep}}\) is spanned by elements of norm at most \(e^{\Lambda_j}\).
\end{itemize}
\emph{Status}: we have read the statements of (D2)--(D5), not their proofs; in particular we do not know which maps \(\varphi_h\) are used.

\begin{problem}\label{pr:low-assembly}
Assume (D1)--(D5), Theorem~\ref{fl-thm:flatat} and Theorem~\ref{thm:uniform}.  Then Lemma~\ref{sg-lem:descent} with \(t = -T_j - (C_F + C_U)u - 2b\), \(P = b\), \(\Lambda = \Lambda_j\) and \(\epsilon = b/\Lambda_j\) gives, if \(C_F \ge 0\) (which we may assume),
\[
  \mu_{\mathrm{core}} \ge -T_j - \bigl(2(C_F + C_U) + 9\bigr)u .
\]
\end{problem}

\begin{theorem}[lower bound]\label{thm:low}
For every admissible family there is \(C_L\) such that, for \(x\) large, at every point and for every lift \(j\), \(\mu_{\mathrm{core}} \ge -T_j - C_Lu\).
\end{theorem}
\noindent\emph{Status}: proved in the formalization.  For \(j < m\) it is Problem~\ref{pr:low-assembly}, from the inputs listed in this chapter and in Theorem~\ref{fl-thm:flatat}; for \(j = m\) it is a separate argument that we have not read.  With Theorem~\ref{sg-thm:reduction} it gives the covolume estimate of Theorem~\ref{hp-in:A4}.

\begin{leanremark}
\raggedright In \cite{Lean} Theorem~\ref{thm:low} is \tl{LowAt}, proved by \tl{low\_\allowbreak{}of\_\allowbreak{}cover\_\allowbreak{}hypotheses}; for \(j < m\) the step of Problem~\ref{pr:low-assembly} is \tl{lowAt\_\allowbreak{}of\_\allowbreak{}cover}, and for \(j = m\) the input is \tl{covering\_\allowbreak{}low\_\allowbreak{}defect\_\allowbreak{}alt}.  Good covers are \tl{Deep} with \tl{Deep.\allowbreak{}Good}, constructed by \tl{exists\_\allowbreak{}good\_\allowbreak{}sized}.  (D1)--(D5) are \tl{SizeAt}, \tl{RingAt}, \tl{MapsAt}, \tl{RkAt}, \tl{GminAt}.  Theorem~\ref{thm:uniform} is \tl{UAt}, from \tl{u\_\allowbreak{}low\_\allowbreak{}of\_\allowbreak{}cover\_\allowbreak{}hyps} and the numerical lemma \tl{UniformLowerBoundOfSums.\allowbreak{}core}; (U1)--(U5) are \tl{SumAt}, \tl{ShareAt}, \tl{DegAt}, \tl{HeightAt}, \tl{HSAt}, and~\eqref{eq:hs} is proved by \tl{hilbertSamuel\_\allowbreak{}of\_\allowbreak{}component\_\allowbreak{}estimates\_\allowbreak{}alt}.  (U4) comes from \tl{height\_\allowbreak{}lower\_\allowbreak{}bound\_\allowbreak{}general\_\allowbreak{}alt} and the statement on covers \tl{DeepStR}, proved by \tl{deepStR\_\allowbreak{}of\_\allowbreak{}use}.  Moriwaki's theorem is \tl{MoriwakiBigStatement}, proved by \tl{moriwakiBig\_\allowbreak{}outright}; the Hodge--Tate theorem is \tl{FarguesHodgeTate}, derived by \tl{farguesHodgeTate\_\allowbreak{}of\_\allowbreak{}farguesExact}.
\end{leanremark}

\section*{Hints}

\subsection*{\S\ref{sec:hs}}
\begin{description}
\item[\ref{pr:binom}] \(\binom{T-1}{\nu} = \prod_{i=1}^\nu(T - i)/\nu!\).
\item[\ref{pr:hs-leading}] Problem~\ref{pr:binom} and the definition of \(h_Z\).
\item[\ref{pr:hs-P1}] Formula~\eqref{cv-eq:hPN} for \(N = 1\); compare with Problem~\ref{cv-pr:degED}.
\item[\ref{pr:log-term}] Minimize over \(y \ge 0\), separately for \(A \le 4\) and \(A > 4\).
\item[\ref{pr:abs-trick}] Problem~\ref{pr:log-term} with \(y = \lvert h\rvert\).
\end{description}

\subsection*{\S\ref{sec:cover}}
\begin{description}
\item[\ref{pr:sum-orth}] Problem~\ref{sm-pr:additivity}.
\item[\ref{pr:weighted-mean}] Bound \(\sum_Z\delta_Z(h_Z + c_s + c_c)\) below by \(\delta'(1-\varrho)\sum_Z\max(0, s_Z + c_s)\), and \(\sum_Z\delta_Z\) above by (U3).
\item[\ref{pr:core}] Problem~\ref{pr:abs-trick} on each component, then Problem~\ref{pr:weighted-mean}; a negative numerator only gets smaller in absolute value when the denominator \(\sum r_Z\) grows.
\end{description}

\subsection*{\S\ref{sec:heights}}
\begin{description}
\item[\ref{pr:rescale-height}] The degree of a metrized line on \(\Oc_{F'}\) changes by \(-[F':\Q]c\).
\item[\ref{pr:small-section}] Compute \(\deg P^*L^{\otimes k}\) with the section \(s(P)\).
\item[\ref{pr:P1-points}] Use the section \(x\) or \(y\) that does not vanish at the point.
\item[\ref{pr:grassmann}] Hadamard's inequality for \(I_a + WW^*\).
\item[\ref{pr:HT-mult}] \(\alpha_G(x)\) is the pullback of \(dT/T\) along \(x \colon G^\vee \to \mu_{p^\infty}\).
\end{description}

\subsection*{\S\ref{sec:descend}}
\begin{description}
\item[\ref{pr:low-assembly}] Theorems~\ref{fl-thm:flatat} and~\ref{thm:uniform} bound \(\mus\) for \(\lambda_-\).
\end{description}

\begin{thebibliography}{Lean}
\bibitem[ALB]{ALB} J.~Ansch\"utz and A.-C.~Le~Bras, \emph{Prismatic Dieudonn\'e theory}, Forum Math.\ Pi \textbf{11} (2023), \url{https://doi.org/10.1017/fmp.2022.22}.
\bibitem[Fa]{Fa} L.~Fargues, \emph{La filtration canonique des points de torsion des groupes \(p\)-divisibles}, Ann.\ Sci.\ \'Ec.\ Norm.\ Sup\'er.\ (4) \textbf{44} (2011), 905--961.
\bibitem[Lean]{Lean} The Lean community, \emph{riemann-hypothesis}, a Lean formalization of Alp\"oge's argument, \url{https://github.com/leanprover-community/riemann-hypothesis}.
\bibitem[Mo]{Mo} A.~Moriwaki, \emph{Continuity of volumes on arithmetic varieties}, J.\ Algebraic Geom.\ \textbf{18} (2009), 407--457.
\bibitem[Yu]{Yu} X.~Yuan, \emph{Big line bundles over arithmetic varieties}, Invent.\ Math.\ \textbf{173} (2008), 603--649.
\bibitem[Zh]{Zh} R.~Zhang, \emph{A preliminary note on Alp\"oge's proposed proof of the generalized Riemann hypothesis}, 22 November 2026, \url{https://hexagonmath.org/pdf/2611.00809v1}.
\bibitem[Zha]{Zha} S.~Zhang, \emph{Positive line bundles on arithmetic varieties}, J.\ Amer.\ Math.\ Soc.\ \textbf{8} (1995), 187--221.
\end{thebibliography}
\end{document}
